\section{Method}

\subsection{Overview}

The wrapper \texttt{FlashAttention3} replaces a pretrained attention module and 
provides three execution paths:

\begin{enumerate}
\item \textbf{FlashAttention-3 kernel} (Hopper + library installed).
\item \textbf{PyTorch SDPA} (any CUDA GPU).
\item \textbf{Head mask + SDPA} (for head pruning experiments).

\end{enumerate}
\textbf{Figure 4.1} shows the dispatch logic.

\begin{lstlisting}
    +----------------------------------+
    |  Input X                          |
    +--------------+-------------------+
                   |
                    
    +----------------------------------+
    |  Q/K/V projection (frozen + LoRA)|
    +--------------+-------------------+
                   |
                    
    +----------------------------------+
    |  Apply RoPE (if applicable)       |
    |  Apply GQA repeat\_kv (if needed)  |
    +--------------+-------------------+
                   |
                    
    +----------------------------------+
    |  Is FA3 available AND CUDA?       |
    +------+-----------------+----------+
           | Yes             | No
                              
    +-------------+   +-----------------+
    | flash\_attn  |   | SDPA fallback   |
    | \_func call  |   | (F.scaled\_...)  |
    +------+------+   +--------+--------+
           |                   |
           +---------+---------+
                      
              +-------------+
              | out\_proj    |
              | (frozen+LoRA)|
              +------+------+
                     |
                      
                    Y
\end{lstlisting}

\subsection{Phase 1: Weight Extraction}

For each projection $P \in \{Q, K, V, O\}$, extract $W_P$ from the source:

\begin{itemize}
\item \textbf{DistilBERT:} \texttt{q\_lin}, \texttt{k\_lin}, \texttt{v\_lin}, \texttt{out\_lin}.
\item \textbf{RoBERTa:} \texttt{self.query}, \texttt{self.key}, \texttt{self.value}; \texttt{output.dense} for \texttt{o}.
\item \textbf{GPT-2:} Split \texttt{c\_attn.weight.T}; \texttt{c\_proj.weight.T} for \texttt{o}.
\item \textbf{Qwen:} \texttt{q\_proj}, \texttt{k\_proj}, \texttt{v\_proj}, \texttt{o\_proj}.

\end{itemize}
Frobenius norms verified $\geq 10^{-4}$.

\subsection{Phase 2: Frozen Base + LoRA}

Each projection becomes a \texttt{LoRALinear}:

\[
\text{Proj}\_P(X) = X W\_P^\top + \frac{\alpha}{r\_L} (X A\_P^\top) B\_P^\top + b\_P
\tag{4.1}
\]

\begin{itemize}
\item $W_P$: frozen (copied from source).
\item $(A_P, B_P)$: LoRA factors, $B = 0$ at init.

\end{itemize}
\textbf{Step-0 preservation.} With $B = 0$, the initial output equals the base 
output.

\subsection{Phase 3: Position and Head Expansion}

\subsubsection{RoPE}

Applied after projection:

\[
Q, K \leftarrow \text{RoPE}(Q, K, \text{cos}, \text{sin})
\tag{4.2}
\]

Supports partial RoPE (rotary\_dim \textless{} head\_dim) via coordinate slicing.

\subsubsection{GQA}

If $H_{kv} < H$:

\[
K \leftarrow \text{repeat\_kv}(K, H/H\_{kv}), \quad
V \leftarrow \text{repeat\_kv}(V, H/H\_{kv})
\tag{4.3}
\]

\subsection{Phase 4: Attention Kernel Dispatch}

\subsubsection{Path A: FlashAttention-3}

If \texttt{FLASH\_ATTN\_AVAILABLE} and \texttt{hidden\_states.is\_cuda}:

\begin{lstlisting}
q\_fa = q.transpose(1, 2).contiguous()  # (B, L, H, d\_h)
k\_fa = k.transpose(1, 2).contiguous()
v\_fa = v.transpose(1, 2).contiguous()

if use\_fp8 and hasattr(torch, "float8\_e4m3fn"):
    q\_fa = q\_fa.to(torch.float8\_e4m3fn)
    k\_fa = k\_fa.to(torch.float8\_e4m3fn)
    v\_fa = v\_fa.to(torch.float8\_e4m3fn)

out = flash\_attn\_func(q\_fa, k\_fa, v\_fa, causal=causal\_flag).to(dtype)
\end{lstlisting}

\textbf{FP8.} If \texttt{use\_fp8=True} and Hopper, casts Q/K/V to FP8.

\subsubsection{Path B: SDPA Fallback}

If FlashAttention-3 fails or is unavailable:

\begin{lstlisting}
if causal\_flag and attn\_mask is None:
    out\_sdpa = F.scaled\_dot\_product\_attention(q, k, v, is\_causal=True)
else:
    out\_sdpa = F.scaled\_dot\_product\_attention(q, k, v, attn\_mask=attn\_mask)
\end{lstlisting}

\textbf{Advantage.} SDPA is always available on PyTorch 2.0+ with CUDA.

\subsection{Phase 5: Output Projection}

\[
Y = \text{Attn}(Q, K, V) W\_O^\top + b\_O
\tag{4.4}
\]

\textbf{Nested bypass.} For RoBERTa/BERT, \texttt{out\_proj} is not applied; the outer 
\texttt{BertSelfOutput.dense} performs it.

\subsection{The Full Algorithm}

\textbf{Algorithm 1:} FlashAttention-3 wrapper forward
\begin{lstlisting}
Input:  X $\in$ R^{B$\times$L$\times$d}, attention\_mask, position\_embeddings
Output: Y, None

1:  Q $\leftarrow$ LoRALinear\_Q(X)
2:  K $\leftarrow$ LoRALinear\_K(X)
3:  V $\leftarrow$ LoRALinear\_V(X)
4:
5:  if RoPE: (Q, K) $\leftarrow$ apply\_rope(Q, K, cos, sin)
6:  if GQA: K, V $\leftarrow$ repeat\_kv(K, V, H/H\_kv)
7:
8:  causal\_flag $\leftarrow$ is\_decoder AND not cross\_attention AND L == kv\_len
9:
10: if FA3\_AVAILABLE and is\_cuda:
11:     try:
12:         if use\_fp8 and fp8\_supported:
13:             Q, K, V $\leftarrow$ cast\_to\_fp8(Q, K, V)
14:         Y\_ctx $\leftarrow$ flash\_attn\_func(Q, K, V, causal=causal\_flag)
15:     except:
16:         Y\_ctx $\leftarrow$ None
17: else:
18:     Y\_ctx $\leftarrow$ None
19:
20: if Y\_ctx is None:                                  Fallback
21:     attn\_mask $\leftarrow$ build\_mask(attention\_mask, causal\_flag)
22:     Y\_ctx $\leftarrow$ SDPA(Q, K, V, attn\_mask)
23:
24: Y\_ctx $\leftarrow$ Y\_ctx.reshape(B, L, H d\_h)
25: if nested: Y $\leftarrow$ Y\_ctx
26: else: Y $\leftarrow$ LoRALinear\_O(Y\_ctx)
27:
28: return Y
\end{lstlisting}

\subsection{Complexity Analysis}

\begin{table}[H]
\centering
\small
\begin{tabular}{lll}
\toprule
\textbf{Path} & \textbf{Time} & \textbf{Memory} \\
\midrule
FA3 & $\mathcal{O}(L^2 d_h)$ + HBM savings & $\mathcal{O}(L d_h)$ \\
SDPA (mem-efficient) & $\mathcal{O}(L^2 d_h)$ & $\mathcal{O}(L d_h)$ \\
SDPA (math) & $\mathcal{O}(L^2 d_h)$ & $\mathcal{O}(L^2)$ \\
\bottomrule
\end{tabular}
\end{table}

\textbf{Actual time savings} depend on hardware. On Hopper, FA3 gives 2-4$\times$ over 
naive; on Ampere/Ada, SDPA memory-efficient path gives 1.5-2$\times$.

\subsection{Parameter Analysis}

\begin{table}[H]
\centering
\small
\begin{tabular}{lll}
\toprule
\textbf{Component} & \textbf{Baseline} & \textbf{Wrapper} \\
\midrule
Base weights & $4 d^2$ & $4 d^2$ (frozen) \\
LoRA ($r_L$) & 0 & $8 r_L d$ \\
\textbf{Total per layer} & $4 d^2$ & $4 d^2 + 8 r_L d$ \\
\bottomrule
\end{tabular}
\end{table}

For DistilBERT ($d=768$, $r_L=16$): \textbf{+589,824 params} across 6 layers.

\textbf{Wrapper increases parameter count by ~66.4M\%.}

\subsection{Theoretical Notes}

\subsubsection{Exactness}

\textbf{Claim.} FlashAttention produces the same output as naive attention up to 
floating-point precision.

\textbf{Proof.} FlashAttention uses online softmax, which is mathematically 
equivalent to standard softmax:

\[
\text{softmax}(s)\_i = \frac{e^{s\_i}}{\sum\_j e^{s\_j}} = \frac{e^{s\_i - m}}{\sum\_j e^{s\_j - m}}
\tag{4.5}
\]

for any shift $m$. The online variant updates $m$ incrementally and rescales 
previous accumulations, giving the same final result. $\square$

\subsubsection{FP8 Accuracy Loss}

\textbf{Claim.} FP8 attention introduces small numerical error.

\textbf{Analysis.} FP8 e4m3 has 3 mantissa bits (~6\% precision). Casting Q, K, V 
to FP8 introduces error $\epsilon \sim 0.06$ per element, which can propagate 
through softmax. Empirically, we observe a small accuracy drop (see Section 8).

\subsection{Hyperparameter Guidance}

\begin{lstlisting}

---
\end{lstlisting}